\subsubsection{História completa}

\paragraph{Muito antes da festa}

\Character{vampirehead} foi uma simples humana, \timestamp{há milhares de anos}. Um dia, enquanto olhava para a lua e para as estrelas e sonhava com a liberdade, isso mudou. Liberdade da fome; liberdade para percorrer os céus e ir a qualquer lado; liberdade da dor de lidar com outros humanos. \Character{vampirehead} desejou desligar-se da humanidade, e algo a ouviu. A partir desse momento, tornou-se vampira: livre de tudo, mas odiada por tudo. O mundo odiava-a tão ferozmente que o próprio sol ameaçava a sua existência.

\Character{vampirehead} foi ficando cada vez mais só, até decidir criar uma família.

Um dia, \timestamp{há três mil anos}, \Character{vampirehead} conheceu \Character{usurper}, uma humana muito interessada em magia e com uma forte sede de poder. \Character{usurper} concebeu uma poção feita com sangue de vampiro que, dada a um humano, o transforma num membro da família. Enquanto dadora do sangue, \Character{vampirehead} podia acrescentar condições ao contrato, mas isso não parecia interessar-lhe.

Tudo pareceu correr bem durante \timestamp{mil anos}, com todas as coisas que duas vampiras podiam fazer para se divertir. Mas a sede de poder de \Character{usurper} desagradava a \Character{vampirehead}. Parecia uma boa familiar, mas era sempre óbvio que a sua prioridade eram as artes místicas.

Isso tornou-se claro à medida que a investigação tomou conta da vida de \Character{usurper}. Havia sempre desarrumação em casa, artefactos espalhados e coisas a desaparecer. \Character{usurper} gostava de que o caos dificultasse a vigilância da sua investigação por \Character{vampirehead}, já que a chefe detestava a desordem. Chegou a cifrar a maior parte das informações nas suas notas.

\timestamp{Há dois mil anos}, \Character{vampirehead} acolheu \Character{enforcer} na família para fazer cumprir as suas regras. A sua palavra só seria superada pela de \Character{vampirehead} na aplicação das regras que esta criou ali mesmo, escrevendo-as num livro: o \Clue{rulebook}.

O livro adquiriu propriedades mágicas: se uma votação da família concluir que uma regra foi violada, o vampiro responsável perde automaticamente os poderes. Isso mostrou a \Character{usurper} como uma vampira original podia ser poderosa, capaz de criar artefactos mágicos quase por acidente.

Sem querer, pelas regras escritas, \Character{vampirehead} estabeleceu que, se perdesse os poderes, \Character{usurper} os herdaria. \Character{usurper} reparou nisso.

\Character{enforcer} memorizava facilmente todas as regras, pelo que escondeu o \Clue{rulebook} no castelo. As regras eram públicas, e no futuro repetiria algumas nos encontros de família.

\timestamp{Décadas depois da morte de Cristo}, a família precisava de melhor gestão doméstica e financeira, por isso convidou \Character{buttler}, uma comerciante conhecida tanto pela competência como pela ganância. Fez bem o seu trabalho. \Character{enforcer} comentava frequentemente que queria uma mordoma só para si.

\timestamp{No ano 100 d.C.}, \Character{enforcer} conheceu um bobo encantador durante uma viagem. Era uma alma artística e carismática, por quem se apaixonou. Como as regras proibiam relações entre vampiros e humanos, pediu a \Character{vampirehead} que transformasse o bobo, \Character{ghoul}, num vampiro.

Mas \Character{enforcer} queria mais: alguém que certamente seguisse as suas regras. Por isso, pediu a \Character{vampirehead} que fizesse a poção com o sangue dela, \Character{enforcer}. Assim, podia acrescentar uma condição ao contrato: ``\Character{ghoul} seguirá sempre as instruções de \Character{enforcer} quando forem dadas como ordens.'' A transformação foi feita em segredo, só com \Character{enforcer} e \Character{ghoul} presentes. \Character{vampirehead} não queria que os outros vampiros tivessem a mesma ideia.

Devido ao seu cargo, \Character{enforcer} precisa de saber o que se passa com todos os outros vampiros, por isso dá a seguinte ordem a \Character{ghoul}: ``Ordeno-te que me contes todas as informações relevantes que descobrires sobre os membros da nossa família, sempre que as descobrires, para todo o sempre.''

Na falsa cerimónia pública de transformação, \Character{buttler} reparou que o frasco dado a \Character{ghoul} não continha a poção. Quando falou disso em privado com \Character{vampirehead}, esta pagou-lhe pelo silêncio.

Com o passar dos anos, juntaram-se vários novos membros à família.

A maioria conspirou, traiu os valores de \Character{vampirehead} e acabou exilada ou morta. Isso mostrou a \Character{usurper} que destronar \Character{vampirehead} não era tarefa simples.

\Character{usurper} ficava responsável por executar os traidores especialmente cruéis quando \Character{enforcer} não estava disponível. Ao fazer experiências durante essas execuções, descobriu que dar a poção proibida a um vampiro o matava.

\Character{supremacist} era uma ditadora terrível que tratava bem os vampiros. A razão era estar a apaixonar-se por \Character{vampirehead}. Juntou-se à família \timestamp{no ano 500 d.C.}. \Character{vampirehead} receava admitir humanos inadequados e, como \Character{supremacist} era tão hostil para com eles, deu-lhe a responsabilidade de decidir em última instância se um humano podia ou não entrar na família.

\Character{lover} era um guerreiro forte e imponente, que serviria de protetor da família. Entrou \timestamp{no ano 700 d.C.}.

\Character{prevlover} era um investigador interessante. Entrou \timestamp{no ano 1000 d.C.}. \Character{usurper} ofereceu-lhe um artefacto que o ajudaria a resolver crimes, o \Clue{clue:truthseeing}. \Character{prevlover} e \Character{enforcer} trabalharam muitas vezes juntos.

\Character{lover} e \Character{traitor} estavam apaixonados, por isso ela juntou-se à família, embora a decisão fosse controversa. Afinal, era uma boa samaritana que andava por aí a fazer boas ações. A ideia era que a imortalidade lhe permitiria fazer ainda mais bem, embora isso fosse, obviamente, um exagero. \Character{traitor} entrou \timestamp{no ano 1400 d.C.}.

A certa altura, \Character{traitor} começou a opor-se à natureza da sua nova família e a conspirar com a \outsideentity{igreja}. Para ela, os vampiros deviam ser curados ou, melhor ainda, destruídos, enquanto os humanos, sobretudo os pobres vindos de outras terras, deviam ser acolhidos.

Toda a gente conhece as opiniões radicais de \Character{traitor} e a considera hipócrita. \Character{usurper} chega a dizer, em tom de ``brincadeira'', que um vampiro morre se beber a poção proibida uma segunda vez e que \Character{traitor} o devia fazer. É um deslize provocado pelo álcool: na verdade, \Character{usurper} receia que \Character{traitor} mate \Character{vampirehead} antes de ela lhe poder herdar os poderes.

Isto contrasta fortemente com \Character{supremacist}, que acredita cada vez mais que todos os humanos devem ser mortos ou transformados em gado. Detesta especialmente que humanos pobres e sem autoridade entrem em países que não são os seus. Para ela, os bons tempos eram aqueles em que governavam os reis.

Ao longo dos anos, \Character{usurper} oferece ajuda a qualquer vampiro que precise e ganha fama de ser a querida ajudante.

\Character{ghoul} é visto como uma piada em vez de um contador de piadas e torna-se ressentido. Quer liberdade e deixa de amar \Character{enforcer}. Na verdade, anseia pela morte dela, pois isso poria fim à obrigação de obedecer às suas ordens.

\Character{supremacist} nunca teve coragem de se declarar a \Character{vampirehead}.

\paragraph{Durante os últimos dois anos}

\timestamp{Há dois anos}, \Character{newcomer} e \Character{vampirehead} conheceram-se. \Character{vampirehead} sentia-se cada vez mais só e convidou-o para uma festa, onde ele conheceu os outros vampiros pela primeira vez.

Todos suspeitavam de que \Character{vampirehead} estava a cortejar \Character{newcomer} e de que seria apenas uma questão de tempo até ele se juntar à família. Provavelmente, só não estavam já juntos porque era contra as regras.

\Character{supremacist} não percebia por que razão \Character{newcomer} devia entrar. Tinha ciúmes da atenção que recebia. \Character{vampirehead} estava tão determinada a acolhê-lo que chegou a ameaçar retirar os poderes a \Character{supremacist}.

Mas, em segredo, \Character{newcomer} e \Character{prevlover} começaram uma relação amorosa.

\timestamp{Há um ano}, \Character{newcomer} e \Character{lover} também começaram uma relação. \Character{newcomer} deixou então \Character{prevlover}, dizendo: ``O que estamos a fazer está errado.''

Enquanto estes casos se desenrolam, \Character{vampirehead} continua sem saber de nada e esforça-se por manter a distância de \Character{newcomer} para não violar as regras.

\Character{traitor} desconfia de \Character{lover} e receia que saiba dos seus planos. Ainda o ama.

\Character{ghoul} rouba uma página do \Clue{rulebook}, dando origem à pista \Clue{clue:hierarchy_page}. Fá-lo porque tem vergonha de ser escravo de \Character{enforcer} e odeia ser constantemente lembrado disso.

Quando \Character{enforcer} repara no desaparecimento, inicia uma investigação. Como \Character{traitor} foi o último membro a entrar e quem mais odeia a família, torna-se a principal suspeita. \Character{enforcer} pede ajuda a \Character{prevlover}. Intercetam mensagens invulgares enviadas de e para a casa de \Character{lover} e de \Character{traitor}. Surgem assim as pistas \Clue{clue:intercepted_messages_loc} e \Clue{clue:intercepted_messages_text}. \Character{prevlover} fica com as notas. Estas mostram que \Character{traitor} conspira com a igreja e que \Character{lover} troca mensagens amorosas com \Character{newcomer}. Contudo, \Character{prevlover} e \Character{enforcer} não dedicam tempo a interpretá-las e não chegam a essa conclusão.

\timestamp{Há dois meses}, como era costume, \Character{vampirehead} marcou \timestamp{16/7/2026} como data da festa para acolher \Character{newcomer} na família. Todos os vampiros tinham de comparecer.

\timestamp{Umas duas semanas antes da festa}, \Character{usurper} perguntou a \Character{vampirehead} se queria ajuda nos preparativos. A chefe agradeceu a oferta e perguntou se mais algum vampiro, além de \Character{usurper}, de \Character{buttler} e de \Character{newcomer}, queria vir na véspera ajudar. Isso significava chegar ao castelo duas noites antes da festa. Todos recusaram e ficaram a saber quem estaria no castelo durante esses dias.

\paragraph{Dois dias antes da festa}

\timestamp{Duas noites antes da festa}, \Character{usurper} e \Character{newcomer} foram para o castelo. \Character{vampirehead} e \Character{buttler} já lá estavam.

\paragraph{Na véspera da festa}

\timestamp{À 01:00}, vários vampiros entram no gabinete de \Character{usurper}.

\Character{ghoul} vai à procura de uma forma de se libertar de \Character{enforcer}. Encontra e rouba um ritual chamado ``Mudar o futuro'', que mais tarde dará origem à pista \Clue{clue:futritualleftovers}.

\Character{prevlover} vai à procura de uma forma de voltar a estar com \Character{newcomer}. Encontra e rouba um ritual chamado ``Mudar e conhecer uma pessoa'', que mais tarde dará origem à pista \Clue{clue:memritualleftovers}. \Character{ghoul} e \Character{prevlover} veem-se. Na pressa, \Character{prevlover} deixa cair a mala, que contém o artefacto que o ajuda a resolver mistérios, o \Clue{clue:truthseeing}, bem como as pistas \Clue{clue:intercepted_messages_loc} e \Clue{clue:intercepted_messages_text}.

\Character{traitor} procura a poção proibida, mas não a encontra. Por isso, encontra uma alternativa: um ritual chamado ``Fazer com que o que é deles seja de outra pessoa'', que mais tarde dará origem à pista \Clue{clue:stealritualleftovers}. Presume que serve para roubar algo e que o poderá usar mais tarde para roubar a poção proibida e executar o seu plano. Na verdade, este ritual desloca todos os objetos do alvo.

\Character{supremacist} rouba um veneno que só funciona em humanos, a \Clue{clue:truth_potion}, ou seja, que mata humanos, mas não vampiros. Na verdade, é uma poção da verdade para vampiros: dá ao vampiro que a beba o poder de exigir respostas verdadeiras a outros seres conscientes.

\timestamp{Às 08:00}, \Character{supremacist} envenena o vinho, dando origem à pista \Clue{clue:poisoned_wine}. Ao sair, encontra \Character{buttler} e discursa longamente sobre como \Character{newcomer} afinal não é assim tão mau e talvez seja uma boa aquisição para a família. Diz que quer que ele beba um vinho melhor do que o que \Character{vampirehead} costuma oferecer, pelo que deixou um vinho especial que só \Character{newcomer} deverá beber. Diz a \Character{buttler} que não quer que se saiba da sua visita nem da mudança de opinião e paga-lhe pelo silêncio sobre a visita e o vinho. Como só precisava de metade do veneno e poderia querer usá-lo mais tarde, guardou a outra metade. Isto dá origem à pista \Clue{clue:truth_potion}.

\Character{vampirehead}, \Character{newcomer}, \Character{buttler} e \Character{usurper} almoçam juntos às \timestamp{12:00}. Quando \Character{buttler} ia servir o vinho, \Character{usurper} impediu-a e disse que tinha trazido o seu próprio vinho de qualidade. Este continha a poção proibida.

\timestamp{Às 15:00}, \Character{ghoul} realiza o ritual e vê o futuro. Vê que \Character{newcomer} vai morrer, embora isso, por si só, não lhe cause grande preocupação. Isto dá origem à pista \Clue{clue:futritualleftovers}. Pela ordem de \Character{enforcer}, tem de revelar tudo o que sabe sobre a família. \Character{ghoul} não percebe logo, mas isto demonstra que \Character{newcomer} já foi transformado.

\Character{enforcer} quer evitar problemas, mas não quer contar a \Character{vampirehead} que andam a roubar coisas a \Character{usurper}, já que \Character{ghoul} está envolvido.

\timestamp{Às 16:00}, \Character{enforcer} telefona a \Character{prevlover} e diz que ignorará o roubo se ele entregar o ritual. Quer usá-lo em \Character{newcomer} para procurar pistas sobre a razão da sua futura morte. Isso permitir-lhe-ia descobrir a relação entre \Character{prevlover} e \Character{newcomer}, o que levaria à perda dos poderes de \Character{prevlover}.

\timestamp{Às 16:15}, \Character{prevlover} entra em pânico e realiza o ritual na memória de \Character{enforcer}. Isto dá origem à pista \Clue{clue:memritualleftovers}.

O resultado imprevisto é apagar memórias aleatórias de todos os vampiros. A maioria diz respeito ao gabinete de \Character{usurper} e a outros artefactos, como as regras do \Clue{rulebook}, que são esquecidas.

\timestamp{Às 17:00}, \Character{traitor} realiza o ritual. Como já não sabe bem o que quer, provoca a dispersão de objetos do gabinete de \Character{usurper} pelo castelo. Esta é a explicação narrativa para a localização da maioria das pistas. Isto dá origem à pista \Clue{clue:stealritualleftovers}.

\timestamp{Às 21:00}, \Character{lover} telefona a \Character{newcomer}. Este sai, e \Character{vampirehead} repara. \Character{traitor} também repara que \Character{lover} não está em casa.

\Character{newcomer} diz a \Character{lover} que a relação terá de acabar, porque \Character{vampirehead} quererá manter a sua própria relação com ele. Talvez um dia, no futuro. \Character{newcomer} e \Character{lover} passam uma última noite juntos.

\paragraph{No dia da festa}

\timestamp{Às 08:00}, \Character{vampirehead} e \Character{newcomer} discutem. Ela quer que ele diga onde esteve na noite anterior, mas ele recusa.

\timestamp{Às 11:00}, \Character{lover} e \Character{newcomer} conversam lá fora. \Character{lover} não quer que a relação termine assim. \Character{buttler} e \Character{usurper} reparam ambas na conversa, sem se aperceberem uma da outra.

Quando a conversa acaba, sem que mais ninguém veja, \Character{buttler} oferece a \Character{lover} a possibilidade de guardar silêncio sobre o assunto, pelo preço certo. \Character{lover} aceita e paga.

\paragraph{Durante a festa}

\timestamp{Durante a transformação}, \Character{newcomer} vai para um quarto com \Character{vampirehead}. Ela dá-lhe a poção proibida correta, convencida de que o vai transformar num vampiro. Deixa-o a descansar e diz aos outros que a poção foi administrada.

Mas \Character{newcomer} já era vampiro. Ao almoço do dia anterior, \Character{usurper} tinha posto a primeira dose apenas no vinho servido a ele. Nem ele nem \Character{vampirehead} o sabiam. A segunda dose mata-o quase de imediato.

\timestamp{Algum tempo depois}, a família vai ver \Character{newcomer} e encontra-o morto. \Character{vampirehead} foi a última pessoa a estar com ele e admite ter-lhe dado a poção. Parece que a chefe matou o mais recente membro da sua própria família.

Era este o plano de \Character{usurper}. Se a família considerar \Character{vampirehead} culpada de homicídio, o \Clue{rulebook} retirar-lhe-á os poderes, e \Character{usurper} herdá-los-á. \Character{vampirehead} deu a dose fatal sem saber o que iria acontecer. \Character{usurper} planeou o homicídio.
